% titel: Blatt 10, A5: Tröte (Stickstoff, einseitig geschlossen)
% skript: Kap. Stehende Wellen in Luftsäulen, Schallgeschwindigkeit im Gas

% Einseitig geschlossenes Rohr (Grundwelle l=lam/4), Stickstoff (kappa=7/5),
% c=sqrt(kappa R T/M), Helium-Verstellung (kappa=5/3, M klein).

Eine Tröte ist ein einseitig \emph{geschlossenes} Rohr: Luft steht an der
geschlossenen Wand still, ist aber am offenen Ende frei. Diese
Randbedingungen bestimmen alle erlaubten Schwingungen.

\begin{werkzeugkasten}[Warum die Grundwelle $\lambda/4$ ist]
An einem \emph{geschlossenen} Ende muss ein Verschiebungs- \emph{Knoten}
sein (die Luft kann dort nicht schwingen, die Wand h"alt sie fest), am
\emph{offenen} Ende ein Bauch (Druckursprung, Luft frei). Das kleinste
Schwingungsmuster, das beides verbindet, ist genau \emph{ein Viertel einer
stehenden Welle}: ein Knoten am geschlossenen, ein Bauch am offenen Ende.
Damit:
\begin{equation}
  L=\frac{\lambda_1}{4}\quad\Rightarrow\quad
  f_1=\frac{c}{\lambda_1}=\frac{c}{4L}\;.
\end{equation}
Die weitere Harmonischenfolge eines einseitig geschlossenen Rohres ist
\textbf{nur ungerade Vielfache} der Grundfrequenz:
\begin{equation}
  f_n=(2n-1)\,f_1,\qquad n=1,2,3,\dots
  \quad(f_1,\,3f_1,\,5f_1,\dots).
\end{equation}
Eine stehende Welle der Länge $(2n-1)\frac{\lambda}{4}$ passt immer genau in
ein solches Rohr, eine gerade Anzahl Viertelwellen nicht (w"urde am offenen
Ende ein Knoten erzwingen). Das unterscheidet das halb offene Rohr von der
beidseitig eingespannten Saite.
\end{werkzeugkasten}

\begin{werkzeugkasten}[Schallgeschwindigkeit im idealen Gas]
Die Schallgeschwindigkeit h"angt von \emph{Steifigkeit} und \emph{Tr"agheit}
des Gases ab, analog zu $c=\sqrt{F/\mu}$ beim Seil:
\begin{equation}
  c=\sqrt{\frac{\kappa R T}{M}},\qquad
  \kappa=\frac{c_p}{c_v}
  \begin{cases}
    7/5 &\text{zweiatomig (N$_2$, O$_2$, Luft)},\\
    5/3 &\text{einatomig (He)},\\
  \end{cases}
\end{equation}
mit der molaren Masse $M$ und der universellen Gaskonstante
$R=\qty{8.314}{J/(mol\,K)}$. Physikalisch: $M$ (Trägheit pro Mol) bremst, die
Temperatur (kinetische Energie) treibt, und $\kappa>1$ wegen der adiabatischen
Verdichtung (die Kompression erw"armt das Gas, "ubertragen auf den
Zusammenhang $\Delta p/\Delta\rho$).

Wichtig: $c$ und damit $f_1=c/4L$ h"angen nicht vom \emph{Druck} ab (bei
idealem Gas), sondern nur von \emph{Temperatur} und molarer Masse. Die
``Normalbedingungen'' $T=\qty{0}{\celsius}=273.15$~K sind hier die relevante
Angabe; der Druck $p=1$~atm fällt aus der Formel heraus, weil $c$ drucklos
ist.
\end{werkzeugkasten}

\begin{werkzeugkasten}[Stickstoff vs.\ Helium -- Vorhersage ohne Rechnung]
Ersetzt man N$_2$ ($\kappa=7/5$, $M=28$~g/mol) durch Helium
($\kappa=5/3$, $M=4$~g/mol), so "andert sich $c\propto\sqrt{\kappa/M}$:
  \begin{itemize}
  \item kleine molare Masse ($28\to4$, Faktor 7) $\Rightarrow$ gr"o"sseres $c$,
  \item gr"o"sseres $\kappa$ ($5/3$ vs.\ $7/5$; einematomig statt zweiatomig)
        $\Rightarrow$ gr"o"sseres $c$.
  \end{itemize}
Beide Effekte wirken in dieselbe Richtung: Die Schallgeschwindigkeit und die
Grundfrequenz der Tröte werden \emph{h"oher} bei Helium. Genauer
$\sqrt{(\kappa_{\mathrm{He}}/\kappa_{\mathrm{N}_2})\cdot(M_{\mathrm{N}_2}/
M_{\mathrm{He}})}=\sqrt{(5/3)/(7/5)\cdot28/4}=\sqrt{25/3}\approx2.89$.
\end{werkzeugkasten}

\begin{achtung}[Einheiten und Größenordnungen]
$c$ in \unit{\metre\per\second}; $RT$ hat die Einheit
$\mathrm{J/mol}=(\mathrm{kg\,m^2/s^2})/\mathrm{mol}$; durch $M$ in
\unit{kg/mol} ergibt sich $\mathrm{m^2/s^2}$, also $c$ in \unit{\metre\per\second}. 
Bei Normalbedingungen $c_{\mathrm{N}_2}\approx\qty{337}{m/s}$, Helium
$\approx\qty{973}{m/s}$ (sehr viel schneller, weil leichtes einatomiges Gas).
Die Länge $l$ der Tröte ist \emph{nicht} numerisch gegeben --- die
Grundfrequenz bleibt "uber $l$ ausgedrückt ($f_1=c/4l$), auch wenn $c$
ausgerechnet wird.
\end{achtung}
